\section{Reproduction Details}
\label{sec:appendix}

\subsection{HapBench Case-level Outcomes}
\label{subsec:app_hapbench_errors}

Table~\ref{tab:hapbench_errors_app} lists every error produced by either executable tool under the fixed 67-case oracle. Cases absent from the table are classified correctly by both tools. We retain the source-annotated oracle even where a more restrictive event or object-lifetime semantics would suggest a different label.

\begin{table}[H]
\centering
\caption{Case-level error inventory for the main HapBench comparison.}
\label{tab:hapbench_errors_app}
\begin{tabularx}{\columnwidth}{@{}l l l Y@{}}
\toprule
Tool & Outcome & Case & Observed mechanism \\
\midrule
HapFlow & FN & \texttt{AnonymousClass2} & Anonymous instance initializer omitted \\
HapFlow & FN & \texttt{Exceptions4} & Exceptional successor/handler not recovered \\
HapFlow & FP & \texttt{ArrayIndexNoLeak} & Distinct constant indices collapse \\
HapFlow & FP & \texttt{VirtualDispatch2} & Infeasible receiver target retained \\
HapFlow & FP & \texttt{VirtualDispatch3} & Infeasible receiver target retained \\
HapFlow & FP & \texttt{UnreachableFlow} & Lifecycle cycle creates reachability \\
\mytool & FN & \texttt{ActivityLifecycle4} & Parent source method is bypassed by child override \\
\mytool & FN & \texttt{BackupExtensionAbility} & Oracle assumes state across extension instances \\
\mytool & FN & \texttt{Button1} & Oracle assumes click-written static state reaches later creation \\
\bottomrule
\end{tabularx}
\end{table}

The two tools agree on 58 outcomes. Of the nine discordant cases, \mytool alone is correct on six and HapFlow alone on three. The exact two-sided McNemar test is therefore based on the discordant pair $(6,3)$ and yields $p=0.508$.

At the finer endpoint-pair unit, \mytool's 52 raw paths canonicalize to 51 source--sink pairs. Source annotations or direct source/IR review confirm 50; the additional pair is the \texttt{NoLeak.log} target at line 14 of \texttt{VirtualDispatch1}, whose allocation receives an empty constant. The retained audit records every pair, its source identity, endpoint lines, provenance, review mode, and decision; an integrity test reproduces TP/FP/FN$=50/1/3$ and 98.04\% precision.

\subsection{Lifecycle-Oracle Sensitivity}
\label{subsec:app_oracle_sensitivity}

The primary evaluation preserves all published/source-annotated HapBench labels. We additionally define a strict executed-order contract for the three disputed lifecycle cases: an override does not execute its parent body without an explicit \texttt{super} call; an instance field does not cross framework-created instances without a same-object witness; and a UI callback occurring after initial ability creation does not feed an earlier \texttt{onCreate} without an explicit later-creation edge. This contract relabels \texttt{ActivityLifecycle4}, \texttt{BackupExtensionAbility}, and \texttt{Button1} as negative and changes no other case.

\begin{table}[H]
\centering
\caption{Sensitivity to strict lifecycle execution order.}
\label{tab:hapbench_strict_oracle}
\resizebox{\columnwidth}{!}{%
\begin{tabular}{@{}lrrrrrrrrr@{}}
\toprule
Tool & TP & TN & FP & FN & Precision & Recall & Specificity & F1 & MCC \\
\midrule
HapFlow~\cite{hapflow} & 48 & 10 & 7 & 2 & 87.27\% & 96.00\% & 58.82\% & 91.43\% & 0.622 \\
\mytool & 50 & 17 & 0 & 0 & 100.00\% & 100.00\% & 100.00\% & 100.00\% & 1.000 \\
\bottomrule
\end{tabular}}
\end{table}

The tools then disagree on nine cases, all classified correctly only by \mytool; the exact two-sided McNemar test gives $p=0.00390625$. This secondary sensitivity contract does not replace the published oracle; it makes the object and event-order semantics responsible for the difference explicit. The artifact derives both tables from the same case predictions through \texttt{analyze\_hapbench\_oracle\_sensitivity.js}.

\subsection{Uncertainty and Category Robustness}

Table~\ref{tab:hapbench_ci_app} reports Wilson intervals for the main configurations. The wide specificity intervals are a direct consequence of having only 14 negative cases; balanced accuracy and MCC are included to make this class imbalance visible.

\begin{table}[H]
\centering
\caption{HapBench uncertainty and category robustness.}
\label{tab:hapbench_ci_app}
\resizebox{\columnwidth}{!}{%
\begin{tabular}{@{}lccccc@{}}
\toprule
Tool & Precision 95\% CI & Recall 95\% CI & Specificity 95\% CI & Macro F1 & Worst-category F1 \\
\midrule
HapFlow~\cite{hapflow} & [82.7, 97.1] & [87.2, 99.0] & [45.4, 88.3] & 95.73\% & 88.89\% \\
\mytool & [92.9, 100.0] & [84.6, 98.1] & [78.5, 100.0] & 97.96\% & 85.71\% \\
$-$IR recovery & [90.6, 100.0] & [56.5, 80.5] & [78.5, 100.0] & 87.07\% & 66.67\% \\
$-$receiver refinement & [89.7, 99.7] & [84.6, 98.1] & [68.5, 98.7] & 97.57\% & 85.71\% \\
Unbounded lifecycle & [90.1, 99.7] & [90.1, 99.7] & [68.5, 98.7] & 98.81\% & 91.67\% \\
\bottomrule
\end{tabular}}
\end{table}

\begin{table}[H]
\centering
\caption{Per-category F1 on all HapBench cases.}
\label{tab:hapbench_categories_app}
\begin{tabular}{@{}lrrr@{}}
\toprule
Category & Cases & HapFlow & \mytool \\
\midrule
Aliasing & 6 & 100.00\% & 100.00\% \\
Anonymous constructs & 10 & 93.33\% & 100.00\% \\
Array-like structures & 5 & 88.89\% & 100.00\% \\
Field/object sensitivity & 6 & 100.00\% & 100.00\% \\
General language features & 22 & 91.89\% & 100.00\% \\
Lifecycle modeling & 13 & 96.00\% & 85.71\% \\
OpenHarmony-specific APIs & 5 & 100.00\% & 100.00\% \\
\bottomrule
\end{tabular}
\end{table}

\subsection{Source-first Benchmark Protocol}
\label{subsec:app_manual_audit}

ArkSourceFirst60 is selected before the evaluated detector run. Projects are partitioned by source-tree size and by whether their imports are compatible with the fixed sensitive-API rules; deterministic hash ordering selects ten projects from each of the six strata. The manifest stores the selection seed, source-tree fingerprint, rule hash, and package-migration hash. Candidate discovery collects the project-wide SDK-import universe and enumerates statically named calls and property reads compatible with that universe, including static element access, named-import calls, and simple callable aliases. It neither reads a \mytool report nor assigns a gold label.

Manual review accepts an occurrence only when its executable source context establishes the configured package/namespace/member and access kind through an import, SDK receiver type, constructor, manager/factory definition, or executable property read. It rejects comments, strings, declarations, same-name application methods, and unrelated library/controller receivers. The delivered gold file contains 90 accepted occurrences; a companion decision file preserves all 186 judgments, including 96 receiver- or owner-incompatible negatives. An integrity test rehashes every source file, validates every exact line and column, checks each canonical rule key, and rejects duplicate occurrence identities. Evaluation joins detector and gold records on project, normalized file, exact line and column, canonical API, and access kind.

\subsection{High-output Stress-audit Protocol}

The 120-project stress audit labels project--namespace--member keys. A key is accepted only when executable source establishes the configured API through an imported namespace/package alias, a manager/helper definition chain, an SDK receiver type, or an executable property rule. Comments, strings, declarations, unrelated same-name methods, and ambiguous method-only matches are rejected. Each accepted key retains project, normalized package/namespace/member, access kind, source file/line, and a bounded source excerpt.

The construction set contains 1,533 reviewed location rows that canonicalize to 666 confirmed project--API keys. The companion integrity audit resolves every key against the delivered source tree and final rules, verifies file/line bounds and snippets, and checks package/namespace/member compatibility. All 666 records pass. Because the set begins from an earlier output ranking, it is not an exhaustive occurrence oracle.

The final detector-only run uses the same frozen build as the corpus and recovers all 666 confirmed keys in all 120 projects. It also reports 182 keys outside the frozen review set. The evaluator therefore defines \emph{confirmed-key reproduction coverage} (666/666) and leaves final-output precision undefined; no unreviewed key enters either side of a precision ratio.

\subsection{Real-project Path-audit Protocol}
\label{subsec:app_path_audit}

The path audit begins from the complete 449-path corpus artifact. A deterministic sampler balances source kind, provenance, sink family, path-length bin, and project diversity, producing 120 records from 48 projects. The fixed queue contains 90 privacy-data and 30 framework-input paths; 64 IFDS, 30 supplementary, and 26 dual-provenance paths; and all available network and UI strata. Each record stores the run-manifest hash, report hash, source/sink file hashes, endpoint locations, path statements, source identity, and provenance.

Reviewers inspect the referenced source and IR path and assign five Boolean decisions: source identity, sink identity, explicit value dependence, reachability consistency, and provenance label. The evaluator accepts a complete path only when all five hold, rejects missing decisions, and rehashes all 240 endpoint files and 120 reports before computing metrics. The audit confirms source identity, sink identity, and provenance for 120/120 paths; explicit dependence, reachability, and complete-path correctness hold for 114/120. All 90 privacy-data paths are complete. The six rejected framework-input paths consist of four unrelated static/component-field joins and two lifecycle/event joins without explicit value dependence.

\subsection{Artifact-to-Paper Traceability}
\label{subsec:app_traceability}

The artifact provides ArkSourceFirst60's selection manifest, complete 186-item decision file, occurrence gold set, candidate-enumerator tests, and integrity evaluator; the 64-case ArkIdentityBench oracle; machine-readable oracles and raw outputs for all 67 HapBench and 48 ArkAsyncBench cases; the Top-120 reproduction evaluator; and the complete semantic-path audit. Evaluators derive confusion matrices, construct/category metrics, endpoint-pair outcomes, paired tests, uncertainty intervals, and runtimes. Figure~\ref{fig:eval_hapbench} and Figure~\ref{fig:eval_corpus_landscape} are generated from evaluator JSON rather than transcribed values.

Corpus tables and figures consume one verified report per project from a single frozen build, SDK, rule set, and command contract; change-impact overlays are rejected by the release summarizer. The manifest records project inventory, build/source/runner hashes, SDK and rule fingerprints, resource bounds, and completion state. The corpus audit checks that every configured path has resolvable source and sink endpoints, a valid \texttt{ifds}, \texttt{async\_supplement}, or \texttt{both} provenance tag, and valid detector-to-path links when the two evidence universes can be joined. Companion JSON retains every plotted value and input manifest hash.
